\lampiransection{Template IFAS, EFAS, dan Matriks SWOT}
\label{app:template-ifas-efas-swot}

Lampiran ini digunakan sebagai lembar kerja awal sebelum hasil akhirnya disajikan pada Bab IV. Faktor awal berikut diturunkan dari observasi GrabFood dan masih perlu diperkuat melalui wawancara, survei, dokumentasi kinerja, dan triangulasi waktu sebelum diberi bobot final.

\begin{table}[!htbp]
    \centering
    \caption{Lembar Kerja Awal Faktor Internal}
    \label{tab:lembar-kerja-faktor-internal}
    \scriptsize
    \setlength{\tabcolsep}{3pt}
    \renewcommand{\arraystretch}{0.92}
    \begin{tabularx}{\textwidth}{|Z{1cm}|Y|Y|Z{2cm}|}
        \toprule
        \textbf{Kode} & \textbf{Faktor Internal Awal} & \textbf{Bukti Data} & \textbf{Kategori} \\ \hline
        \midrule
        S1 & Reputasi digital cukup kuat karena rating 4,7 didukung 4.584 penilaian. & TO-GF-01; Lampiran~\ref{app:evidence-observasi-grabfood-20260613} & Kekuatan \\ \hline
        S2 & Rasa, porsi, bahan segar, kemasan, variasi menu, dan pelanggan rutin muncul sebagai tema positif. & TO-GF-02; TO-GF-07 & Kekuatan \\ \hline
        S3 & Menu terlaris dan promo/voucher aktif memberi bahan untuk positioning dan promosi penjualan. & TO-GF-03; TO-GF-04 & Kekuatan \\ \hline
        W1 & Ulasan negatif/netral menunjukkan risiko salah item, item kurang lengkap, add-on tidak diberikan, dan koordinasi layanan. & TO-GF-08 & Kelemahan \\ \hline
        W2 & Menu atau produk tidak tersedia dapat mengganggu persepsi kelengkapan pilihan dan pengalaman pemesanan. & TO-GF-06 & Kelemahan \\ \hline
        W3 & Rentang harga yang lebar perlu dijelaskan dengan nilai porsi/lauk agar tidak dibaca sebagai harga mahal. & TO-GF-05; TO-KP-06 & Kelemahan potensial \\ \hline
        \bottomrule
    \end{tabularx}
    \tablesource{Diolah peneliti sebagai lembar kerja IFAS awal berdasarkan observasi GrabFood (2026).}
\end{table}

\begin{table}[!htbp]
    \centering
    \caption{Lembar Kerja Awal Faktor Eksternal}
    \label{tab:lembar-kerja-faktor-eksternal}
    \scriptsize
    \setlength{\tabcolsep}{3pt}
    \renewcommand{\arraystretch}{0.92}
    \begin{tabularx}{\textwidth}{|Z{1cm}|Y|Y|Z{2cm}|}
        \toprule
        \textbf{Kode} & \textbf{Faktor Eksternal Awal} & \textbf{Bukti Data} & \textbf{Kategori} \\ \hline
        \midrule
        O1 & Promo/voucher GrabFood dapat dimanfaatkan untuk menarik pelanggan baru dan mendorong pembelian ulang. & TO-GF-03 & Peluang \\ \hline
        O2 & Tema positif ulasan dapat diperkuat sebagai bukti sosial pada komunikasi produk dan tampilan toko. & TO-GF-02; TO-GF-07 & Peluang \\ \hline
        O3 & Kompetitor juga memiliki catatan perbaikan menu, akurasi pesanan, rasa, porsi, dan pemenuhan permintaan sehingga konsistensi layanan dapat menjadi pembeda. & TO-KP-05 & Peluang \\ \hline
        T1 & Pondok Gurih memiliki rating 4,9, 56.157 penilaian, serta penanda \textit{Legendaris} dan Resto Bintang 5. & TO-KP-01; TO-KP-07 & Ancaman reputasi \\ \hline
        T2 & Paripurna berjarak 3,0 km dari titik observasi, sedangkan Pondok Gurih 5,0 km dan Garuda 6,7 km tetap berada dalam ruang pembanding akses pelanggan. & TO-KP-03 & Ancaman akses \\ \hline
        T3 & Dendeng Batokok memiliki batas harga atas lebih rendah, sementara Istana Krakatau, Pondok Gurih, dan Garuda memiliki rentang premium yang mendekati atau melampaui harga atas Nasi Gerilya. & TO-KP-06 & Ancaman harga dan nilai \\ \hline
        T4 & Garuda dan Istana Krakatau memiliki akumulasi penilaian lebih besar daripada Nasi Gerilya walaupun ratingnya lebih rendah. & TO-KP-01; TO-KP-02 & Ancaman bukti sosial \\ \hline
        T5 & Ulasan platform berisiko bias sehingga persepsi pelanggan yang tidak menulis review belum tentu terlihat. & TO-GF-07; TO-GF-08; Lampiran~\ref{app:template-triangulasi-data} & Ancaman metodologis \\ \hline
        \bottomrule
    \end{tabularx}
    \tablesource{Diolah peneliti sebagai lembar kerja EFAS awal berdasarkan observasi GrabFood dan kompetitor (2026).}
\end{table}

\begin{table}[!htbp]
    \centering
    \caption{Template Penilaian IFAS}
    \label{tab:template-penilaian-ifas}
    \footnotesize
    \setlength{\tabcolsep}{4pt}
    \renewcommand{\arraystretch}{0.95}
    \begin{tabularx}{\textwidth}{|Z{1.2cm}|Y|c|c|c|}
        \toprule
        \textbf{Kode} & \textbf{Faktor Internal} & \textbf{Bobot} & \textbf{Rating} & \textbf{Skor} \\ \hline
        \midrule
        S1 & Reputasi digital dan jumlah penilaian & 0,00 & 0 & 0,00 \\ \hline
        S2 & Kekuatan rasa, porsi, kemasan, dan pelanggan rutin & 0,00 & 0 & 0,00 \\ \hline
        S3 & Menu terlaris dan promo aktif & 0,00 & 0 & 0,00 \\ \hline
        W1 & Akurasi dan kelengkapan pesanan & 0,00 & 0 & 0,00 \\ \hline
        W2 & Ketersediaan menu dan pembaruan stok & 0,00 & 0 & 0,00 \\ \hline
        W3 & Persepsi harga dan nilai porsi/lauk & 0,00 & 0 & 0,00 \\ \hline
        \textbf{Total} &  & \textbf{1,00} &  & \textbf{0,00} \\ \hline
        \bottomrule
    \end{tabularx}
    \tablesource{Diolah peneliti sebagai template IFAS; bobot dan rating diisi setelah triangulasi (2026).}
\end{table}

\begin{table}[!htbp]
    \centering
    \caption{Template Penilaian EFAS}
    \label{tab:template-penilaian-efas}
    \footnotesize
    \setlength{\tabcolsep}{4pt}
    \renewcommand{\arraystretch}{0.95}
    \begin{tabularx}{\textwidth}{|Z{1.2cm}|Y|c|c|c|}
        \toprule
        \textbf{Kode} & \textbf{Faktor Eksternal} & \textbf{Bobot} & \textbf{Rating} & \textbf{Skor} \\ \hline
        \midrule
        O1 & Pemanfaatan promo/voucher GrabFood & 0,00 & 0 & 0,00 \\ \hline
        O2 & Penguatan bukti sosial dari ulasan positif & 0,00 & 0 & 0,00 \\ \hline
        O3 & Diferensiasi melalui konsistensi layanan dibanding kompetitor & 0,00 & 0 & 0,00 \\ \hline
        T1 & Reputasi kompetitor sangat kuat & 0,00 & 0 & 0,00 \\ \hline
        T2 & Kompetitor lebih dekat atau sebanding dari sisi akses & 0,00 & 0 & 0,00 \\ \hline
        T3 & Kompetitor memiliki rentang harga bawah dan premium yang kuat & 0,00 & 0 & 0,00 \\ \hline
        T4 & Akumulasi penilaian kompetitor tertentu lebih besar & 0,00 & 0 & 0,00 \\ \hline
        T5 & Bias ulasan platform dalam membaca pengalaman pelanggan & 0,00 & 0 & 0,00 \\ \hline
        \textbf{Total} &  & \textbf{1,00} &  & \textbf{0,00} \\ \hline
        \bottomrule
    \end{tabularx}
    \tablesource{Diolah peneliti sebagai template EFAS; bobot dan rating diisi setelah triangulasi (2026).}
\end{table}

\begin{table}[!htbp]
    \centering
    \caption{Template Alternatif Strategi SWOT Berdasarkan Temuan Awal}
    \label{tab:template-alternatif-strategi-swot}
    \scriptsize
    \setlength{\tabcolsep}{3pt}
    \renewcommand{\arraystretch}{0.92}
    \begin{tabularx}{\textwidth}{|Z{2.4cm}|Y|Y|}
        \toprule
        \textbf{Faktor} & \textbf{Kekuatan (S)} & \textbf{Kelemahan (W)} \\ \hline
        \midrule
        \textbf{Peluang (O)} & Strategi SO: gunakan rating, jumlah penilaian, menu terlaris, ulasan positif, dan promo aktif untuk memperkuat paket unggulan serta komunikasi nilai porsi/rasa di GrabFood. & Strategi WO: gunakan promo, survei pelanggan, dan pembaruan tampilan menu untuk memperbaiki persepsi harga, stok, dan pengalaman pesanan. \\ \hline
        \textbf{Ancaman (T)} & Strategi ST: gunakan bukti sosial, menu unggulan, konsistensi rasa, dan komunikasi nilai porsi untuk menghadapi Pondok Gurih yang kuat secara reputasi, Garuda/Istana Krakatau yang kuat secara akumulasi penilaian, Paripurna yang lebih dekat, serta Dendeng Batokok yang kuat sebagai pembanding harga bawah. & Strategi WT: perbaiki SOP cek pesanan, pembaruan stok aplikasi, penanganan komplain, dan monitoring ulasan agar kelemahan layanan tidak memperbesar ancaman kompetitor maupun bias pembacaan review. \\ \hline
        \bottomrule
    \end{tabularx}
    \tablesource{Diolah peneliti sebagai template matriks SWOT awal berdasarkan observasi GrabFood (2026).}
\end{table}
